% ============================================================
\section{Experimental Evaluation}
\label{sec:genact-evaluation}
% ============================================================

The objective of this evaluation is to investigate the utility of GenACT as a
configurable data generator for experiments involving evolving knowledge and
stream reasoning. In particular, the experiments examine whether the generated
data can produce workloads with different temporal, structural, and reasoning
characteristics, and whether these differences are observable when the data are
processed by existing reasoning systems.

GenACT is designed around three characteristics of evolving knowledge:
\emph{temporality}, \emph{dynamicity}, and \emph{timeliness}. In addition to
generating temporally ordered RDF data, it supports partitioning according to
domain attributes, such as conferences and users, and according to structural
graph patterns, such as stars, chains, and trees. Furthermore, the ACE ontology
is available in OWL~2 EL, QL, RL, and DL variants. Consequently, GenACT
provides several dimensions along which an experimental reasoning workload can
be configured while retaining a common semantic domain.

The evaluation is guided by the following questions:

\begin{itemize}

    \item[\textbf{Q1.}] Can GenACT generate heterogeneous temporal workloads
    exhibiting different activity patterns and query behaviours?

    \item[\textbf{Q2.}] How does increasing the number of simultaneously active
    conference streams affect continuous query processing?

    \item[\textbf{Q3.}] Do different competency questions expose different
    computational and result characteristics over the generated streams?

    \item[\textbf{Q4.}] How can GenACT expose workload and performance
    variations when evolving stream data are combined with ontology-mediated
    reasoning?

\end{itemize}

To investigate these questions, we evaluate the thirteen competency questions
using C-SPARQL2 and RDFox. The two systems serve complementary purposes:
C-SPARQL2 exercises continuous RDF stream processing, whereas RDFox allows the
generated data to participate in OWL~2 RL reasoning. The objective is
therefore not to compare the performance of the two systems, but to
demonstrate the different reasoning workloads that can be constructed from the
same GenACT data.

This split reflects a genuine gap in the current landscape of RDF stream
processing: continuous-query engines and expressive ontology reasoners are
largely disjoint families of systems, and no single mature implementation
combines both. We therefore considered candidates from each family separately
before settling on C-SPARQL2 and RDFox, and report the outcome for every
system considered rather than silently narrowing the field.

Among \emph{RDF stream processing engines} (no OWL reasoning), three systems
were considered. \textbf{C-SPARQL2} was selected and used throughout this
chapter: it builds and runs reliably, and its continuous-query model matches
GenACT's own window-based competency questions directly.
\textbf{RSP4J}~\cite{rsp4j}, together with its reference engine
\textbf{Yasper} (the two are not separate systems: Yasper is RSP4J's own
built-in implementation of the RSP-QL model it defines), was built from
source and confirmed to execute the same window-close evaluation semantics as
C-SPARQL2 and RDFox on a reduced instance of the real GenACT stream. At the
full evaluation scale used throughout this chapter, however, it did not
complete within a practical time budget: a single competency question against
one conference's real stream (approximately 9{,}400 files) ran for over 80
minutes of continuous, genuine computation -- confirmed via live CPU-time
monitoring, not a deadlock -- without terminating, against seconds for the
same workload under C-SPARQL2 or RDFox. RSP4J/Yasper is therefore not used for
comparative numbers in this chapter; it is documented here as a real,
empirically-grounded scalability limitation on this workload, not a system
that was skipped. \textbf{CQELS}~\cite{cqel} was identified as a candidate but
not attempted directly, for three concrete reasons established while
investigating StreamQR (below): its only available implementation targets the
pre-2015 Jena package namespace rather than the modern one used throughout
this pipeline; RSP4J's own maintainers ship a CQELS integration module
(\texttt{cqels-rsp4j}) but leave it commented out of their default build, a
signal that this integration path is not currently maintained; and CQELS
exposes its own native continuous-query dialect rather than RSP-QL, so a
harness would need a third distinct query interface. Given that RSP4J, judged
the more tractable of the two going in, still consumed a substantial
investigation before hitting a genuine wall, CQELS was assessed as at least as
likely to fail and is left for future work rather than attempted.

Among \emph{stream-enabled DL reasoners}, two systems were considered.
\textbf{RDFox} was selected because it supports materialisation-based
reasoning with real incremental maintenance as data changes, and its sandbox
scripting interface could be driven directly from real GenACT output.
\textbf{StreamQR}~\cite{streamqr} was also considered, and a substantial
build effort was made: its build tooling was repaired across five sequential,
independent issues (dead Scala-plugin repositories, an sbt syntax break, a
JDK-version incompatibility in its Scala compiler bootstrap, and missing
generated parser sources regenerated from its raw grammar), after which
dependency resolution completed successfully. It ultimately could not be
reproduced, however, because two of its core dependencies
(\texttt{rsp:cqels\_2.11:1.1.0} and \texttt{eu.trowl:trowl-core:1.4}) are not
published in any publicly reachable repository and appear to be private build
outputs from the original authors, never released. Real substitutes were
built from source for both, but StreamQR's own code is written against a
newer Jena API than those substitutes expose, so it does not compile against
them. This is a genuine dependency dead end rather than an unattempted
system, and StreamQR is consequently not used for comparative numbers in this
chapter.


% ============================================================
\subsection{Setup}
\label{subsec:genact-experimental-setup}
% ============================================================

We evaluate GenACT using C-SPARQL2 and RDFox to exercise complementary
reasoning settings. C-SPARQL2 evaluates continuous queries over RDF streams,
whereas RDFox allows the evolving ABox to participate in OWL~2 RL
materialisation. Both experiments use the same thirteen competency questions
and $N\in\{1,2,5\}$ concurrently active conference streams. The systems are not
evaluated as direct competitors; rather, they expose different aspects of the
workloads enabled by GenACT.

The experiments were performed on a machine equipped with an Intel Core
i7-14650HX processor with 16 cores and 24 threads, 16~GB of memory, and
Windows~11. The generated RDF data are replayed directly from GenACT's
\texttt{.ttl} files over TCP for C-SPARQL2. Consequently, the experiments
operate on the actual output of the generator rather than on a separately
constructed benchmark dataset. Each competency question is evaluated
independently using a dedicated C-SPARQL2 instance: executing several
competency questions within one engine instance was found to interfere with
their individual evaluation frequencies, so isolating each query lets its
behaviour reflect the corresponding GenACT workload rather than contention
with unrelated queries.

% ------------------------------------------------------------
\subsubsection{Window Configuration}
\label{subsubsec:genact-window-config}
% ------------------------------------------------------------

The thirteen competency questions represent different analytical requirements
and therefore operate over different temporal scopes. Each competency question
is assigned a window range, determining how much recent stream data is
available to the query, and a window step, determining how frequently the
active window advances, as shown in Table~\ref{tab:genact-window-tiers}.

\begin{table}[htbp]
\centering
\caption{Window configurations used for the GenACT competency questions.}
\label{tab:genact-window-tiers}
\begin{tabular}{clrr}
\toprule
\textbf{Tier} & \textbf{Competency Questions} &
\textbf{Range (ms)} & \textbf{Step (ms)} \\
\midrule
Short  & CQ1, CQ7, CQ13
            & 60{,}000  & 2{,}000 \\

Medium & CQ2, CQ3, CQ5, CQ6, CQ9, CQ10, CQ12
            & 60{,}000  & 2{,}000 \\

Long   & CQ4, CQ8, CQ11
            & 180{,}000 & 6{,}000 \\
\bottomrule
\end{tabular}
\end{table}

The same range and step definitions are used when constructing the
corresponding C-SPARQL2 and RDFox workloads. This preserves the temporal
semantics of the competency questions across the two experiments even though
the underlying systems implement the workload differently.

The relationship between evaluation latency and the configured step is
particularly important. A latency value in isolation does not indicate whether
a stream-processing engine is able to keep pace with the requested evaluation
frequency. If the average evaluation time becomes larger than the window step,
the engine requires more time to evaluate a window than the interval at which a
new evaluation is requested.

The competency questions provide a system-independent specification of the
information needs. For the experiments, they are instantiated using the query
constructs supported by C-SPARQL2 and RDFox; the executable query sets are
provided in the accompanying repository.

% ------------------------------------------------------------
\subsubsection{Evaluation Metrics}
\label{subsubsec:genact-metrics}
% ------------------------------------------------------------

For every C-SPARQL2 run, we record four complementary measurements.

\emph{Average latency} denotes the average wall-clock time required by
C-SPARQL2 to evaluate one window. \emph{Throughput} reports the number of
completed window evaluations per second of experiment runtime.
\emph{Average memory} denotes the average JVM heap usage observed during the
run. Finally, \emph{non-empty windows} reports the percentage of window
evaluations that produce at least one result satisfying the complete query
pattern.

The non-empty-window measure provides important context when interpreting the
performance measurements. An empty window is not an erroneous evaluation; it
may simply correspond to a temporal interval in which the required combination
of graph patterns is absent. Hence, reporting query activity together with
latency and throughput helps distinguish computational cost from query
selectivity.

For RDFox, we report query-evaluation latency and the proportion of non-empty
windows for the windowed experiment, and additionally record the number of
asserted and materialised facts in a bulk experiment to verify that the
generated ABox participates in actual OWL~2 RL inference.

For the C-SPARQL2 experiment, we additionally normalize average query latency
by the configured window step:

\[
    \rho =
    \frac{\text{average evaluation latency}}
         {\text{window step}}.
\]

A value $\rho < 1$ indicates that the average evaluation completes before the
next scheduled window evaluation, whereas $\rho > 1$ indicates that the
evaluation requires more time than the interval between successive evaluations.
This measure is specific to the continuous-query execution model of C-SPARQL2
and is therefore not used to compare C-SPARQL2 with RDFox.


% ============================================================
\subsection{Patterns and Peaks}
\label{subsec:genact-patterns-peaks}
% ============================================================

Before evaluating the generated data with reasoning systems, we first examine
its temporal characteristics. A useful stream-reasoning workload should not
only vary in size, but should also exhibit changes in event density over time.

Figure~\ref{fig:patternsnpeaks} illustrates the weekly tweet-volume patterns
generated for five conference streams. Each conference individually exhibits
characteristic activity peaks associated with important phases of the
conference lifecycle. These include, for example, early announcements, review
notifications, and the period during which the conference itself is taking
place.

When the conference streams are merged, however, these peaks occur at
different positions in time and overlap with one another. The resulting stream
therefore exhibits a substantially more heterogeneous distribution of event
volumes than any individual conference stream.

\begin{figure}[htbp]
    \centering
    \includegraphics[width=\textwidth]{figures/owl2stream/peaks.png}
    \caption{Tweet-volume patterns and peaks for five conferences:
    individual conference streams (left) and the merged stream (right).}
    \label{fig:patternsnpeaks}
\end{figure}

This property is important for the subsequent experiments. Increasing the
number of active conferences does not simply replicate a uniform input stream.
Instead, it combines conferences that may be in different stages of their
lifecycle and therefore contribute different quantities of data at different
points in time. As a result, a reasoning system is exposed to fluctuating
workloads and temporally concentrated bursts of information.

Thus, temporal workload intensity represents one configurable dimension of
GenACT, complementing other dimensions such as the number of concurrent
conference streams, the selected competency question, graph structure, and
ontology expressivity.


% ============================================================
\subsection{Evaluation with C-SPARQL2}
\label{subsec:genact-csparql}
% ============================================================

The C-SPARQL2 experiment evaluates all thirteen competency questions under
$N=1$, $N=2$, and $N=5$ simultaneously active conference streams, using the
setup and window configuration described above. Here, $N$ represents genuine
concurrent stream-processing load: the conference streams are active
simultaneously and their continuous queries compete for the available
computational resources.

This experiment therefore primarily addresses the stream-processing
dimensions of the evaluation, including the effects of increasing concurrent
stream volume and differences between competency-question workloads.


% ------------------------------------------------------------
\subsubsection{Results}
\label{subsubsec:genact-csparql-results}
% ------------------------------------------------------------

Table~\ref{tab:genact-csparql-results} reports the complete results for all
thirteen competency questions under $N\in\{1,2,5\}$ parallel conference
streams. The table provides the exact numerical measurements, while
Figure~\ref{fig:genact-csparql-performance} provides a visual summary of
latency, throughput, and memory consumption across the different workloads.

\begin{table*}[htbp]
\centering
\caption{C-SPARQL2 results for all thirteen GenACT competency questions under
$N\in\{1,2,5\}$ parallel conference streams.}
\label{tab:genact-csparql-results}
\scriptsize
\setlength{\tabcolsep}{3.5pt}
\resizebox{\textwidth}{!}{%
\begin{tabular}{cl rrr rrr rrr rrr}
\toprule
&
&
\multicolumn{3}{c}{\textbf{Latency (ms)}} &
\multicolumn{3}{c}{\textbf{Throughput (win/s)}} &
\multicolumn{3}{c}{\textbf{Memory (MB)}} &
\multicolumn{3}{c}{\textbf{Non-empty (\%)}} \\
\cmidrule(lr){3-5}
\cmidrule(lr){6-8}
\cmidrule(lr){9-11}
\cmidrule(lr){12-14}

\# & \textbf{Competency Question}
& $N=1$ & $N=2$ & $N=5$
& $N=1$ & $N=2$ & $N=5$
& $N=1$ & $N=2$ & $N=5$
& $N=1$ & $N=2$ & $N=5$ \\
\midrule

1 & Trending Topics
& 814.3 & 683.4 & 1{,}739.2
& 0.50 & 0.57 & 0.42
& 73.2 & 72.6 & 112.6
& 74.0 & 86.2 & 79.1 \\

2 & Active Research Groups
& 342.4 & 922.3 & 3{,}644.0
& 0.50 & 0.86 & 0.30
& 72.4 & 83.6 & 151.6
& 100.0 & 100.0 & 100.0 \\

3 & Publication Activities
& 880.6 & 976.7 & 3{,}823.1
& 0.51 & 0.86 & 0.28
& 71.0 & 79.6 & 159.6
& 74.5 & 68.2 & 51.4 \\

4 & Conference Match
& 1{,}194.1 & 1{,}671.7 & 5{,}330.7
& 0.17 & 0.31 & 0.20
& 75.4 & 98.0 & 181.4
& 75.0 & 73.3 & 60.0 \\

5 & Interdisciplinary Authors
& 1{,}032.6 & 1{,}212.6 & 5{,}933.2
& 0.48 & 0.84 & 0.18
& 69.0 & 77.4 & 155.8
& 72.9 & 66.7 & 45.5 \\

6 & Session Popularity
& 802.3 & 925.7 & 3{,}355.2
& 0.51 & 0.86 & 0.31
& 70.8 & 80.6 & 157.2
& 24.0 & 11.6 & 20.0 \\

7 & Global Research Focus
& 1{,}000.2 & 1{,}108.6 & 3{,}652.7
& 0.50 & 0.85 & 0.30
& 93.8 & 103.6 & 168.9
& 73.5 & 68.2 & 55.6 \\

8 & Funding Organizations
& 1{,}026.5 & 1{,}535.7 & 3{,}896.3
& 0.18 & 0.32 & 0.28
& 75.1 & 97.7 & 168.7
& 25.0 & 20.0 & 20.6 \\

9 & Networking Opportunities
& 815.5 & 919.1 & 3{,}209.1
& 0.51 & 0.86 & 0.34
& 70.8 & 81.0 & 152.6
& 100.0 & 100.0 & 97.7 \\

10 & Collaboration Networks
& 954.4 & 1{,}050.3 & 2{,}395.8
& 0.49 & 0.84 & 0.46
& 69.9 & 80.1 & 147.6
& 73.5 & 69.4 & 56.0 \\

11 & Non-academic Collaborators
& 3{,}034.7 & 3{,}663.1 & 10{,}814.0
& 0.16 & 0.28 & 0.09
& 98.5 & 114.9 & 272.4
& 76.5 & 66.7 & 68.8 \\

12 & Geographical Distribution
& 831.0 & 1{,}530.7 & 1{,}802.5
& 0.51 & 0.68 & 0.59
& 72.7 & 75.5 & 141.4
& 100.0 & 100.0 & 90.8 \\

13 & Platform Impact
& 796.7 & 1{,}528.4 & 1{,}857.8
& 0.51 & 0.68 & 0.60
& 72.8 & 79.0 & 142.0
& 74.0 & 67.1 & 59.4 \\

\bottomrule
\end{tabular}%
}
\end{table*}


% ------------------------------------------------------------
% PERFORMANCE VISUALISATION
% ------------------------------------------------------------

\begin{figure*}[htbp]
\centering

\begin{tikzpicture}

\begin{groupplot}[
    group style={
        group size=3 by 1,
        horizontal sep=1.2cm
    },
    width=0.34\textwidth,
    height=0.58\textwidth,
    enlargelimits=false,
    axis on top,
    xtick={0,1,2},
    xticklabels={$N=1$,$N=2$,$N=5$},
    ytick={0,1,2,3,4,5,6,7,8,9,10,11,12},
    yticklabels={
        CQ1,CQ2,CQ3,CQ4,CQ5,CQ6,CQ7,
        CQ8,CQ9,CQ10,CQ11,CQ12,CQ13
    },
    y dir=reverse,
    xlabel={Parallel conference streams},
    tick label style={font=\scriptsize},
    label style={font=\small},
    title style={font=\small\bfseries},
    colorbar style={
        tick label style={font=\scriptsize}
    }
]

% LATENCY
\nextgroupplot[
    title={(a) Average latency},
    ylabel={Competency question},
    point meta min=0,
    point meta max=11000,
    colorbar,
    colorbar style={
        ylabel={Latency (ms)}
    }
]

\addplot[
    matrix plot*,
    mesh/cols=3,
    point meta=explicit
]
table[meta=value] {
x y value
0 0 814.3
1 0 683.4
2 0 1739.2
0 1 342.4
1 1 922.3
2 1 3644.0
0 2 880.6
1 2 976.7
2 2 3823.1
0 3 1194.1
1 3 1671.7
2 3 5330.7
0 4 1032.6
1 4 1212.6
2 4 5933.2
0 5 802.3
1 5 925.7
2 5 3355.2
0 6 1000.2
1 6 1108.6
2 6 3652.7
0 7 1026.5
1 7 1535.7
2 7 3896.3
0 8 815.5
1 8 919.1
2 8 3209.1
0 9 954.4
1 9 1050.3
2 9 2395.8
0 10 3034.7
1 10 3663.1
2 10 10814.0
0 11 831.0
1 11 1530.7
2 11 1802.5
0 12 796.7
1 12 1528.4
2 12 1857.8
};

% THROUGHPUT
\nextgroupplot[
    title={(b) Throughput},
    yticklabels=\empty,
    point meta min=0,
    point meta max=0.9,
    colorbar,
    colorbar style={
        ylabel={Windows/s}
    }
]

\addplot[
    matrix plot*,
    mesh/cols=3,
    point meta=explicit
]
table[meta=value] {
x y value
0 0 0.50
1 0 0.57
2 0 0.42
0 1 0.50
1 1 0.86
2 1 0.30
0 2 0.51
1 2 0.86
2 2 0.28
0 3 0.17
1 3 0.31
2 3 0.20
0 4 0.48
1 4 0.84
2 4 0.18
0 5 0.51
1 5 0.86
2 5 0.31
0 6 0.50
1 6 0.85
2 6 0.30
0 7 0.18
1 7 0.32
2 7 0.28
0 8 0.51
1 8 0.86
2 8 0.34
0 9 0.49
1 9 0.84
2 9 0.46
0 10 0.16
1 10 0.28
2 10 0.09
0 11 0.51
1 11 0.68
2 11 0.59
0 12 0.51
1 12 0.68
2 12 0.60
};

% MEMORY
\nextgroupplot[
    title={(c) Average memory},
    yticklabels=\empty,
    point meta min=60,
    point meta max=280,
    colorbar,
    colorbar style={
        ylabel={Memory (MB)}
    }
]

\addplot[
    matrix plot*,
    mesh/cols=3,
    point meta=explicit
]
table[meta=value] {
x y value
0 0 73.2
1 0 72.6
2 0 112.6
0 1 72.4
1 1 83.6
2 1 151.6
0 2 71.0
1 2 79.6
2 2 159.6
0 3 75.4
1 3 98.0
2 3 181.4
0 4 69.0
1 4 77.4
2 4 155.8
0 5 70.8
1 5 80.6
2 5 157.2
0 6 93.8
1 6 103.6
2 6 168.9
0 7 75.1
1 7 97.7
2 7 168.7
0 8 70.8
1 8 81.0
2 8 152.6
0 9 69.9
1 9 80.1
2 9 147.6
0 10 98.5
1 10 114.9
2 10 272.4
0 11 72.7
1 11 75.5
2 11 141.4
0 12 72.8
1 12 79.0
2 12 142.0
};

\end{groupplot}
\end{tikzpicture}

\caption{Visual summary of the C-SPARQL2 performance measurements for the
thirteen GenACT competency-question workloads under one, two, and five
parallel conference streams. The panels show average latency, throughput,
and average JVM memory consumption.}
\label{fig:genact-csparql-performance}

\end{figure*}


% ------------------------------------------------------------
\subsubsection{Effect of Increasing Conference Streams}
\label{subsubsec:genact-scaling}
% ------------------------------------------------------------

Table~\ref{tab:genact-csparql-results} and
Figure~\ref{fig:genact-csparql-performance} show that increasing the number of
simultaneously active conference streams creates a substantially more demanding
workload for C-SPARQL2, particularly for the $N=5$ configuration.

For example, CQ2 (\emph{Active Research Groups}) increases from an average
latency of 342.4~ms for one conference stream to 3,644.0~ms for five streams.
CQ3 (\emph{Publication Activities}) increases from 880.6~ms to 3,823.1~ms,
while CQ5 (\emph{Interdisciplinary Authors}) increases from 1,032.6~ms to
5,933.2~ms.

The most computationally demanding workload is CQ11
(\emph{Non-academic Collaborators}), whose average latency increases from
3,034.7~ms at $N=1$ to 10,814.0~ms at $N=5$.

A similar scaling effect can be observed in memory consumption. For CQ2,
average memory usage increases from 72.4~MB to 151.6~MB, while CQ11 increases
from 98.5~MB to 272.4~MB. The $N=5$ configurations generally require
substantially more memory than their corresponding single-stream workloads.

The results are not strictly monotonic for every metric at every intermediate
configuration. For example, CQ1 has a lower measured latency for $N=2$ than
for $N=1$, and several $N=2$ configurations exhibit higher measured throughput
than the corresponding single-stream run. The relevant observation is therefore
not that every intermediate measurement increases monotonically, but that the
larger $N=5$ configuration produces a clear increase in processing pressure
across the workload as a whole.


% ------------------------------------------------------------
\subsubsection{Workload Heterogeneity}
\label{subsubsec:genact-workload-heterogeneity}
% ------------------------------------------------------------

The experiments also show that the number of streams alone does not determine
the computational characteristics of a GenACT workload. Different competency
questions exhibit markedly different behaviour even when evaluated using the
same number of conference streams.

For $N=1$, CQ2 requires an average of only 342.4~ms per evaluation, whereas
CQ11 requires 3,034.7~ms. At $N=5$, CQ12
(\emph{Geographical Distribution}) requires 1,802.5~ms, while CQ11 requires
10,814.0~ms.

The non-empty-window ratios provide further evidence of workload diversity.
CQ2 produces results in 100\% of its evaluated windows for all three stream
configurations. CQ9 similarly produces results in 100.0\%, 100.0\%, and
97.7\% of its windows. By contrast, CQ6 produces non-empty results in only
24.0\%, 11.6\%, and 20.0\% of the corresponding windows, while CQ8 remains
around 20--25\%.

Importantly, the increase in processing cost cannot therefore be attributed
simply to whether a query produces results. CQ2 remains non-empty in every
window while its average latency increases from 342.4~ms to 3,644.0~ms.
Similarly, CQ9 remains almost completely non-empty while its latency increases
from 815.5~ms to 3,209.1~ms.

These differences arise from the semantic and structural characteristics of the
competency questions. The queries introduce different graph patterns, joins,
selectivities, and interactions with static background knowledge. Thus, the
thirteen competency questions form a heterogeneous workload suite derived from
a common temporal domain rather than thirteen variations of one uniform query.


% ------------------------------------------------------------
\subsubsection{Ability to Maintain the Window Schedule}
\label{subsubsec:genact-window-pressure}
% ------------------------------------------------------------

The relationship between latency and the configured window step provides a
particularly useful interpretation of these results. To make this relationship
explicit, Figure~\ref{fig:genact-step-normalized-latency} reports the average
evaluation latency normalized by the step assigned to each competency question:

\[
\rho =
\frac{\text{average evaluation latency}}
     {\text{window step}}.
\]

A value of $\rho < 1$ means that the average query evaluation completes within
the interval before the next evaluation is requested. A value of $\rho > 1$
indicates that evaluation requires more time than the configured step.

\begin{figure}[htbp]
\centering

\begin{tikzpicture}

\begin{axis}[
    width=\textwidth,
    height=7.2cm,
    ybar,
    bar width=2.3pt,
    ymin=0,
    ylabel={Latency / window step},
    xlabel={Competency question},
    symbolic x coords={
        CQ1,CQ2,CQ3,CQ4,CQ5,CQ6,CQ7,
        CQ8,CQ9,CQ10,CQ11,CQ12,CQ13
    },
    xtick=data,
    x tick label style={
        rotate=45,
        anchor=east,
        font=\scriptsize
    },
    tick label style={font=\scriptsize},
    label style={font=\small},
    legend style={
        at={(0.5,1.12)},
        anchor=south,
        legend columns=3,
        font=\small
    },
    enlarge x limits=0.03
]

% N = 1
\addplot coordinates {
(CQ1,0.407)
(CQ2,0.171)
(CQ3,0.440)
(CQ4,0.199)
(CQ5,0.516)
(CQ6,0.401)
(CQ7,0.500)
(CQ8,0.171)
(CQ9,0.408)
(CQ10,0.477)
(CQ11,0.506)
(CQ12,0.416)
(CQ13,0.398)
};

% N = 2
\addplot coordinates {
(CQ1,0.342)
(CQ2,0.461)
(CQ3,0.488)
(CQ4,0.279)
(CQ5,0.606)
(CQ6,0.463)
(CQ7,0.554)
(CQ8,0.256)
(CQ9,0.460)
(CQ10,0.525)
(CQ11,0.611)
(CQ12,0.765)
(CQ13,0.764)
};

% N = 5
\addplot coordinates {
(CQ1,0.870)
(CQ2,1.822)
(CQ3,1.912)
(CQ4,0.888)
(CQ5,2.967)
(CQ6,1.678)
(CQ7,1.826)
(CQ8,0.649)
(CQ9,1.605)
(CQ10,1.198)
(CQ11,1.802)
(CQ12,0.901)
(CQ13,0.929)
};

% rho = 1 boundary
\addplot[
    domain=-1:13,
    samples=2,
    dashed,
    thick
] {1};

\legend{$N=1$, $N=2$, $N=5$}

\end{axis}

\end{tikzpicture}

\caption{Average C-SPARQL2 evaluation latency normalized by each competency
question's configured window step. Values above $1$ indicate that an average
evaluation requires more time than the interval between two successive window
evaluations.}
\label{fig:genact-step-normalized-latency}

\end{figure}

For $N=1$ and $N=2$, all competency questions remain below the corresponding
window-step threshold. At $N=5$, however, CQ2, CQ3, CQ5, CQ6, CQ7, CQ9,
and CQ10 exceed their 2~s step. CQ11 also exceeds its longer 6~s step,
reaching an average evaluation latency of 10.814~s.

This observation is more informative than an increase in latency alone. It
shows that GenACT can produce workloads that move a stream-processing engine
between different operational regimes: from one in which the requested
continuous-query schedule can be maintained to one in which the engine begins
to fall behind the rate at which evaluations are requested.


% ============================================================
\subsection{Structural Workloads}
\label{subsec:genact-structural-workloads}
% ============================================================

The number of concurrent conferences represents only one way of varying a
GenACT workload. The generated data can additionally be partitioned according
to domain attributes or according to its graph structure.

Attribute-based partitioning can, for example, create individual streams for
conferences, users, organisations, research domains, or other entities present
in the generated data. Structural partitioning instead identifies graph
patterns such as stars, chains, and trees.

These patterns represent different connectivity structures. Star-shaped graphs
contain a central node connected to several related entities, chain structures
capture sequences of relations, and tree structures represent hierarchical
relationships. Such differences may produce different join and reasoning
characteristics even for datasets of similar size.

To validate this functionality, the structural \texttt{CONSTRUCT} queries were
executed against the complete production EventData corpus rather than a
selected sample. The corpus contains 26,834 generated EventData files.

\begin{table}[htbp]
    \centering
    \caption{Structural patterns identified in the generated EventData corpus.}
    \label{tab:genact-shapes}
    \begin{tabular}{lrr}
        \toprule
        \textbf{Shape} &
        \textbf{Matching Files} &
        \textbf{Coverage} \\
        \midrule
        Star  & 26{,}834 / 26{,}834 & 100.0\% \\
        Chain & 24{,}966 / 26{,}834 & 93.0\% \\
        Tree  & 0 / 26{,}834 & 0.0\% \\
        \bottomrule
    \end{tabular}
\end{table}

All files satisfy the star-pattern query, while 24,966 files, corresponding to
approximately 93\% of the corpus, also satisfy the chain pattern. No file
satisfies the tree pattern. This is not a sampling artefact: the hardcoded
generator's per-tweet-template design allows one file to assert several
objects under a single predicate, but no template gives one of those objects a
further branching relationship within the same file, so a genuine tree is
structurally unreachable by this pipeline rather than merely rare.

Structural partitioning therefore provides another workload axis that is
independent of simply increasing the amount of data. A researcher can construct
streams dominated by particular graph topologies and investigate whether a
reasoning system behaves differently for highly centralised, sequential, or
hierarchical input structures.


% ============================================================
\subsection{Evaluation with RDFox}
\label{subsec:genact-rdfox}
% ============================================================

GenACT is also intended for settings in which evolving instance data interact
with expressive ontological knowledge. The C-SPARQL2 experiments demonstrate
the utility of GenACT for continuous RDF stream processing, but they evaluate
graph patterns over explicitly available triples and do not exploit the OWL
axioms represented by the ACE ontology. Since current stream reasoners do not
provide sufficient OWL~2 expressivity to exercise this dimension directly, as
discussed above, we use RDFox with the ACE-RL ontology
as a complementary experimental setting. The objective is to examine how
GenACT exposes workload variation when evolving data participate in
ontology-mediated reasoning, rather than to compare RDFox directly with
C-SPARQL2.

RDFox supports Datalog-based materialisation and can be used with the OWL~2 RL
variant of the ACE ontology. The generated ABox can therefore be loaded
together with the ACE-RL TBox, materialised, and subsequently queried over
knowledge containing both explicitly asserted and inferred facts. A
preliminary execution without an ontology or rules is insufficient for this
purpose: in that configuration RDFox behaves essentially as an RDF store. The
relevant evaluation therefore uses the reasoning capabilities provided by
materialisation over the ACE-RL ontology, which remains loaded while the
active ABox changes with each window.

% ------------------------------------------------------------
\subsubsection{Verification of OWL~2 RL Reasoning}
\label{subsubsec:genact-rdfox-verification}
% ------------------------------------------------------------

Before evaluating the windowed workload, we first verify that loading the
generated ABox together with ACE-RL results in actual inferred knowledge
rather than RDFox operating solely as a triple store. The experimental
configuration is summarized in Table~\ref{tab:genact-rdfox-setup}.

\begin{table}[htbp]
    \centering
    \caption{Configuration of the RDFox evaluation.}
    \label{tab:genact-rdfox-setup}
    \begin{tabular}{ll}
        \toprule
        \textbf{Parameter} & \textbf{Configuration} \\
        \midrule
        Reasoner             & RDFox \\
        Ontology             & ACE OWL~2 RL \\
        Instance data        & GenACT-generated conference data \\
        Reasoning strategy   & Materialisation \\
        Query semantics      & Asserted and entailed knowledge \\
        Conference scale     & $N \in \{1, 2, 5\}$ \\
        \bottomrule
    \end{tabular}
\end{table}

For each scale $N$, the generated ABox was first imported into a fresh,
rule-less RDFox store to obtain a clean asserted-fact count. The ACE-RL TBox
was then imported into the same store, which triggers RDFox's real
incremental Datalog materialisation rather than a plain bulk triple load, and
the number of explicitly asserted facts is compared with the number of facts
available after materialisation. Table~\ref{tab:genact-rdfox-results} reports
the results.

\begin{table}[htbp]
\centering
\caption{RDFox OWL~2 RL materialisation over GenACT-generated data.}
\label{tab:genact-rdfox-results}
\small
\begin{tabular}{crrrr}
\toprule
\textbf{Configuration} &
\textbf{Asserted} &
\textbf{Materialised} &
\textbf{Additional inferred} &
\textbf{Materialisation time (s)} \\
\midrule

$N=1$ & 57{,}382  & 57{,}974  & 592     & 0.071 \\
$N=2$ & 97{,}787  & 98{,}804  & 1{,}017 & 0.095 \\
$N=5$ & 153{,}143 & 154{,}537 & 1{,}394 & 0.084 \\

\bottomrule
\end{tabular}
\end{table}

The materialised knowledge base is larger than the explicitly asserted ABox in
all three configurations. RDFox derives 592 additional facts for $N=1$,
1{,}017 for $N=2$, and 1{,}394 for $N=5$. These differences provide direct
evidence that OWL~2 RL reasoning is active and produces additional knowledge
from the generated data. Materialisation time, meanwhile, stays under 100~ms
throughout and does not grow monotonically with corpus size (0.071~s,
0.095~s, and 0.084~s for $N=1$, $N=2$, and $N=5$, respectively); query
latency against the materialised store is below one millisecond at every
scale, since the store is RAM-resident and remains small enough that
materialisation, not querying, dominates total cost.

% ------------------------------------------------------------
\subsubsection{Windowed Ontology-Mediated Evaluation}
\label{subsubsec:genact-rdfox-windowed}

The bulk result above loads an entire conference's data once and materialises
over it, which is not directly comparable to the C-SPARQL2 evaluation's
repeated per-window evaluation. To obtain a genuine window-for-window
comparison, all thirteen competency questions were additionally evaluated
against RDFox using the same sliding-window definitions as
Table~\ref{tab:genact-window-tiers} (same range and step per tier, and the
same replay-rate scaling used to pace the C-SPARQL2 streams), but computed
directly from the real timestamps embedded in the generated files rather than
from a live TCP replay. This was necessary because the original paper's own
RDFox streaming harness -- built on an embedded JNI client library from RDFox
7.1a -- no longer matches the API or native library of the RDFox~7.6
installation used throughout this thesis, and porting it surfaced a further,
unrelated defect: the shared TCP replay/receiver code (used by both engines)
only flushes accumulated window data from inside its read loop, so any data
received after the last scheduled flush is silently never committed once the
stream ends. Rather than continue patching 2015-era code three layers deep,
a new, independent pipeline was built directly on the RDFox~7.6 installation:
for each window, files whose embedded timestamp falls in that window's range
are incrementally added to the store, files that have fallen out of range are
incrementally deleted, and the query is re-evaluated -- with the ACE-RL TBox
and, where relevant, the static organisation background data loaded once at
the start of each run, so every window benefits from real, incrementally
maintained OWL~2~RL materialisation rather than a cold reload. As with
C-SPARQL2, each of $N=1$, 2, and 5 conference streams is evaluated
independently and the results aggregated, rather than merging conferences
onto one timeline; unlike an earlier pass of this experiment, the $N$
conference streams now run as $N$ genuinely concurrent RDFox processes
(one real OS process per conference, launched at the same time and awaited
together), matching how C-SPARQL2's own $N$ parallel engine instances
actually compete for the same CPU and memory, rather than one conference
being evaluated after another.

\begin{table*}[htbp]
\centering
\caption{RDFox windowed evaluation (with OWL~2~RL materialisation) for all thirteen GenACT competency questions under $N\in\{1,2,5\}$ parallel conference streams. Latency is real query evaluation time reported by RDFox; \emph{$<$1\,ms} values round to 0.000\,s at RDFox's own reporting precision.}
\label{tab:genact-rdfox-windowed}
\scriptsize
\setlength{\tabcolsep}{4pt}
\resizebox{\textwidth}{!}{%
\begin{tabular}{cl rrr rrr}
\toprule
& & \multicolumn{3}{c}{\textbf{Latency (ms)}} & \multicolumn{3}{c}{\textbf{Non-empty (\%)}} \\
\cmidrule(lr){3-5}
\cmidrule(lr){6-8}
\# & \textbf{Competency Question} & $N=1$ & $N=2$ & $N=5$ & $N=1$ & $N=2$ & $N=5$ \\
\midrule
1  & Trending Topics            & 17.5 & 17.8 & 18.5 & 59.3 & 59.3 & 59.0 \\
2  & Active Research Groups     & 6.6  & 7.2  & 5.5  & 72.9 & 72.9 & 70.0 \\
3  & Publication Activities     & $<$1 & $<$1 & $<$1 & 30.5 & 30.5 & 25.7 \\
4  & Conference Match           & 0.8  & 0.7  & 0.8  & 60.0 & 60.0 & 59.6 \\
5  & Interdisciplinary Authors  & 0.8  & 0.6  & 0.6  & 57.6 & 56.8 & 54.7 \\
6  & Session Popularity         & $<$1 & $<$1 & $<$1 & 16.9 & 16.9 & 16.3 \\
7  & Global Research Focus      & $<$1 & $<$1 & $<$1 & 30.5 & 30.5 & 25.7 \\
8  & Funding Organizations      & $<$1 & $<$1 & $<$1 & 25.0 & 25.0 & 23.1 \\
9  & Networking Opportunities   & $<$1 & $<$1 & $<$1 & 72.9 & 72.9 & 70.0 \\
10 & Collaboration Networks     & 0.8  & 0.9  & 0.8  & 57.6 & 57.6 & 55.7 \\
11 & Non-academic Collaborators & $<$1 & $<$1 & $<$1 & 60.0 & 60.0 & 59.6 \\
12 & Geographical Distribution  & 26.1 & 45.9 & 32.1 & 72.9 & 72.9 & 70.0 \\
13 & Platform Impact            & 15.8 & 19.5 & 13.1 & 59.3 & 59.3 & 59.0 \\
\bottomrule
\end{tabular}%
}
\end{table*}

Two things stand out against the C-SPARQL2 table
(Table~\ref{tab:genact-csparql-results}), and both survived moving from
sequential to genuinely concurrent execution of the $N$ conference streams.
First, non-empty rates stay exactly flat as $N$ grows within a stream count
(e.g.\ CQ2: 72.9\,\%~$\to$~72.9\,\%~$\to$~70.0\,\%; CQ9:
72.9\,\%~$\to$~72.9\,\%~$\to$~70.0\,\%), rather than falling the way
C-SPARQL2's do -- expected, since non-empty rate is a property of the window
contents themselves (computed once per conference, independent of how many
other conferences happen to run alongside it), not of execution timing.
Second, and less expected: latency does \emph{not} show the same clear,
monotonic climb with $N$ that C-SPARQL2 exhibits, even now that the $N$
conference streams are real, concurrently-running RDFox processes competing
for the same CPU and memory. A few CQs move in the expected direction (CQ1:
17.5\,ms~$\to$~17.8~$\to$~18.5), but others move the opposite way (CQ2:
6.6~$\to$~7.2~$\to$~5.5\,ms) or non-monotonically (CQ12: 26.1~$\to$~45.9~$\to$~32.1\,ms).
The explanation was checked empirically rather than assumed: monitoring a
single RDFox process's live thread count (via the operating system's process
table, polled every 50\,ms from launch to exit) showed it using up to 29
threads on its own -- already \emph{more} than the 24 logical processors this
machine has, for a single instance. RDFox's own internal parallelism therefore
already oversubscribes this machine's hardware at $N=1$, before any
concurrent conference streams are added. Raising $N$ to 2 or 5 multiplies the
total number of threads contending for the same 24 real ones (up to
$5\times29=145$ at $N=5$), but since the baseline is already oversubscribed,
the added contention is not a clean, proportional function of $N$: which run
happens to get favourable operating-system scheduling on a given pass becomes
close to non-deterministic, which is consistent with the non-monotonic
pattern observed above. This is a real property of running RDFox at this
scale on this hardware, not a flaw in the measurement. Non-empty rates broadly agree with C-SPARQL2's own ranking of
which competency questions are selective (CQ6, CQ8 sparse; CQ2, CQ9 dense),
which is expected since both evaluate the same triple patterns against the
same underlying corpus -- the exact percentages differ because the two
systems compute window membership differently (live TCP arrival time for
C-SPARQL2 versus the files' own embedded timestamps here). A separate check
confirms RL entailment is not a contributing factor here: removing the ACE-RL
TBox entirely and re-running reproduces every one of the 39 non-empty
percentages above (thirteen competency questions times three stream counts)
bit-for-bit, and per-CQ latency stays within the same order of magnitude with
or without the TBox loaded (e.g.\ CQ1: 15.7--18.5\,ms either way; CQ12:
26.6--45.9\,ms either way) -- consistent with the earlier finding that this
corpus is small enough for materialisation cost to be negligible. The gap
against C-SPARQL2 is therefore attributable entirely to the window-computation
difference, not to RDFox answering over additional entailed facts or being
fast because it skips reasoning.

The purpose of the RDFox evaluation is not to establish that one system is
faster than the other -- the two use fundamentally different execution models,
and the windowed comparison above is offered in the same spirit as
C-SPARQL2's own table, not as a claim of superiority. What the RDFox
experiments add that C-SPARQL2 cannot is the ontology-mediated dimension: the
same generated conference data, evaluated under real OWL~2~RL materialisation,
with results reflecting both explicitly asserted and inferred knowledge.

% ============================================================
\subsection{GenACT as a Configurable Workload Generator}
\label{subsec:genact-workload-generator}
% ============================================================

Taken together, the experiments show that GenACT exposes several independent
dimensions for constructing reasoning workloads.

\begin{table}[htbp]
    \centering
    \caption{Principal workload dimensions exposed by GenACT.}
    \label{tab:genact-workload-dimensions}
    \begin{tabular}{p{3.2cm}p{4.0cm}p{6.2cm}}
        \toprule
        \textbf{Dimension} &
        \textbf{Configuration} &
        \textbf{Resulting workload characteristic} \\
        \midrule

        Stream volume &
        Number of simultaneous conferences &
        Amount of temporally active knowledge that must be processed \\

        Temporal behaviour &
        Event lifecycle, peaks, window range, and step &
        Event bursts, temporal density, and required evaluation frequency \\

        Query workload &
        Thirteen competency questions &
        Different joins, selectivities, and result densities \\

        Static/stream interaction &
        Background knowledge combined with generated streams &
        Joins and enrichment across static and evolving knowledge \\

        Entity partitioning &
        Conference, user, organisation, research domain, etc. &
        Different numbers and compositions of independent streams \\

        Graph topology &
        Star, chain, and tree partitions &
        Different connectivity and join structures \\

        Ontology expressivity &
        ACE ontologies for OWL~2 EL, QL, RL, and DL &
        Different ontology-mediated reasoning requirements \\

        \bottomrule
    \end{tabular}
\end{table}

This configurability is more significant than any individual latency
measurement. A workload generator that only increases the number of triples
primarily varies dataset size. GenACT additionally allows the nature of the
workload itself to be changed while preserving a coherent domain model.

For example, an experiment may hold the competency question fixed and increase
the number of concurrently active conference streams. Another may maintain a
similar data scale but select a different competency question and therefore a
different join and selectivity pattern. Structural partitioning may be used to
construct streams with different graph topologies, while ontology profile
selection allows the generated ABox to participate in reasoning tasks with
different levels of expressivity.

GenACT therefore provides a common experimental basis for workload dimensions
that are often studied separately: temporal stream processing, static-background
enrichment, structural graph variation, and ontology-mediated reasoning.


% ============================================================
\subsection{Discussion}
\label{subsec:genact-evaluation-discussion}
% ============================================================

The evaluation demonstrates that GenACT provides both \emph{workload
scalability} and \emph{workload heterogeneity}.

The C-SPARQL2 experiments show that increasing the number of concurrently active
conference streams can substantially increase latency and memory consumption.
For the largest configuration, several competency questions require more time
for an average evaluation than their configured window step, indicating that
the engine can no longer maintain the requested continuous-query schedule.

At the same time, workload characteristics differ considerably between
competency questions even for a fixed number of streams. Some competency
questions produce matching results in almost every window, whereas others are
substantially more selective. Their latency and memory characteristics also
differ markedly. These effects arise from meaningful semantic differences in
the query workload rather than from arbitrary modifications to RDF triples.

The temporal peak analysis provides another dimension. Multiple conference
streams are not uniform copies of one another: conferences can be in different
phases of their lifecycle and therefore contribute different event volumes at
different times. Merging them results in fluctuating stream intensity and
temporally concentrated workloads.

Structural partitioning further demonstrates that the generated corpus need not
be treated as one monolithic stream. Streams may be organised according to
entities or graph topology, creating alternative workloads from the same
underlying generated domain.

Finally, the RDFox evaluation complements the stream-processing experiments by
exercising the ontology-mediated dimension of GenACT. This distinction is
particularly relevant given the present separation between systems optimised
for continuous RDF stream processing and systems supporting more expressive OWL
reasoning. By generating both temporally rich instance data and corresponding
OWL~2 ontologies, GenACT can support experiments on both sides of this
landscape.

Consequently, GenACT should not be viewed solely as a mechanism for producing
synthetic RDF triples. Its contribution lies in providing a configurable
experimental environment in which stream volume, temporal behaviour, query
semantics, graph topology, stream composition, and reasoning expressivity can
be varied while retaining a common and interpretable semantic domain.


% ============================================================
\subsection{Summary}
\label{subsec:genact-evaluation-summary}
% ============================================================

The experimental evaluation highlights four properties of GenACT.

First, the generated workloads are \emph{temporally heterogeneous}. Individual
conference streams exhibit event-driven activity peaks, while combining
multiple conferences produces changing stream volumes and overlapping bursts.

Second, the workloads are \emph{scalable}. Increasing the number of concurrent
conference streams substantially increases the processing burden and, for some
configurations, causes continuous-query evaluation to exceed the configured
window step.

Third, the workloads are \emph{structurally and semantically diverse}. The
thirteen competency questions exercise different joins, selectivities, static
background knowledge, and result densities, while attribute- and shape-based
partitioning provide additional mechanisms for constructing alternative stream
structures.

Finally, the generated data support \emph{different levels of reasoning
expressivity}. C-SPARQL2 demonstrates continuous RDF stream processing over
GenACT workloads, while the RDFox evaluation uses the ACE OWL~2 RL ontology to
exercise ontology-mediated materialisation.

Together, these properties demonstrate the usefulness of GenACT as a
configurable workload generator for studying reasoning over evolving knowledge.
